\section{Lastne oblike in frekvence kvadratne opne}

V tej nalogi obravnavamo problem lastnih vrednosti nihanja opne v 2D
\begin{equation}
    \nabla^2\;u = -\lambda\;u,
\end{equation}
kjer je $\lambda=k^2$ lastna vrednost. Robni pogoji so Dirichletovi $u=0$, valovno število $k^2=\frac{\omega^2\,\rho\,d}{\gamma}$, kjer je $\omega=2\,\pi\,f$ krožna frekvenca, $\rho=1$ gostota, $d=1$ debelina in $\gamma=1$ linijska obremenitev opne na robu.

V Kartezijsvem koordinatnem sistemu se Laplaceov operator diskeritiza po metodi končnih razlik v matriko $\mathbf{A}$, ki ima vrednosti 4 na diagonali, na preostalih štirih izven-diagonalnih delih ima vrednost -1. Matriko delimo še z kvadratom koraka $h=\Delta x=\Delta y$, nakar lahko uporabimo metode za reševanje problema lastnih vrednosti, kot so \\\texttt{scipy.linalg.eig}, \texttt{scipy.sparse.eigs}, Jakobijeva iteracija ... oziroma inverzno potenčno metodo v treh vrsticah
\begin{align}
    y_{i+1} &= (A-\sigma\,I)^{-1}\cdot x_i\\\nonumber
    x_{i+1} &= \frac{y_{i+1}}{\norm{y_{i+1}}}\\\nonumber
    \lambda_{i+1} &= x_{i+1}^T\cdot A\cdot x_{i+1},
\end{align}
če nas zanima le posamezna lastna vrednosti najbližje parametru $\sigma$.

\subsection{Homogena opna}

V tem primeru je matrika $\mathbf{A}$ simetrična, zato lahko uporabimo metode za Hermitske matrike (npr. \texttt{numpy.linalh.eigh}), kar občutno pospeši izračun, obenem pa ni potrebe po uporabi posplošenega problema lastnih vrednosti z matriko desnih strani $B$.

Na sledečih slikah so prikazane lastne oblike in vrednosti za različne metode, kjer opazimo, da so rezultati primerljivi, sploh za prvih 4 lastnih oblik. Pri višjih ima Jacobi težave, prav tako pa rihaja do manjših odstopanj med direktnim reševanjem iz z algoritmom ARPACK. Vseeno je MAC kriterij (Modal Assurance Criterion) v vseh primerih precej visok, zato ga ne prikazujemo.

\fig{H}{1}{./6.nal_slike/homogena/lastne_homog}{Lastne oblike in frekvence homogene opne}{}
\fig{H}{0.7}{./6.nal_slike/homogena/sparse_homog}{ARPACK iz \texttt{scipy.sparse}}{}
\fig{H}{0.7}{./6.nal_slike/homogena/jacobi_homog}{Jacobi algoritem diagonalizacije}{}

Opazimo, da lastne frekvence (in posledično vrednosti) zavzemajo diskretne nivoje, kjer se na nekaterih nahaja tudi več lastnih oblik, katerim je skupno to, da imajo enako število vrhov.

ARPACK algoritem, ki je v ozadju \texttt{scipy.sparse.linalg.eigsh}, omogoča izračun manjšega števila lastnih vrednosti, zato sem analizairal le prvih 9 najmanjših. računski čas je primerljiv z direktno metodo (\texttt{scipy.linalg.eigh}), vendar je pri velikih matrikah ARPACK učinkoviteši, ako potrebujemo le nekaj lastnih oblik. Jacobi je očitno precej počasnejši, reda $\mathcal{O}(n^3)$, čemur verjetno pripomore tudi slabša implementacija v Pythonu.

\fig{H}{0.7}{./6.nal_slike/homogena/frekvence_homog}{Lastne frekvence homogene opne za različne metode}{}
\fig{H}{0.7}{./6.nal_slike/homogena/cas_homog_30000s}{Računski čas -- primerjava metod}{}

%vpeljava gostote!
\subsection{Nehomogena opna}

Če je opna nehomogena, se enačba spremeni
\begin{equation}
    \nabla^2\;u = -\lambda\,\rho\,u,
\end{equation}
kjer gostota $\rho$ ni več konstantna, lastna vrednost pa predstavlja $\lambda = \frac{\omega^2\,d}{\gamma}$. Tako sledi po MKR diskretiziran zapis
\begin{equation}
    \mathbf{A}\,\mathbf{u} = -\lambda\,\mathbf{B}\,\mathbf{u},
\end{equation}
pri katerem gre za posplošeni problem lastnih vrednosti, saj imamo na desni strani matriko $B$, ki je slična gostoti, zaradi česar ostane matrika $A$ simetrična.

Za drug način pa enačbo zapišemo v obliki
\begin{equation}
    \frac{1}{\rho}\;\nabla^2\;u = -\lambda\;u,
\end{equation}
kjer po dikretizaciji z MKR matrike $B$ ni, zato pa $A$ postane nesimetrična! To pomeni, da moramo uporabiti splošne metode za reševanje problema lastnih vrednosti cele matrike, ki so občutno počasnejše.

Na spodnji sliki je prikazan izračun lastnih vrednosti homogene opne po različnih metodah, kjer posplošeno označuje funkcijo iz SciPy \texttt{eigh} oz. \texttt{eigsh} za redke (sparse) matrike in za nesimetrične \texttt{eig} oz. \texttt{eigs}. Pri tem velja poudariti, da je bilo potrebno pri funkcijah za redke matrike lastne vrednosti še dodatno razvrstiti po velikosti.
\fig{H}{0.7}{./6.nal_slike/homogena/eigval_homog}{Lastne vrednosti homogene opne}{}


% \fig{H}{0.8}{./6.nal_slike/}{}{}

\subsubsection{`f'-opna}

Uporabimo porazdelitev gostote kot na predlogi, kjer belo območje predstavlja $\rho_1=1$, črno pa $\rho_2=ratio$. Pri tem $ratio=\frac{\rho_2}{\rho_1}$ predstavlja razmerje gostot.
\fig{H}{0.3}{./6.nal_slike/opna_f}{Porazdelitev gostote}{}

%%Ratio=100
Najprej preverim konvergenco pri $ratio = 100$. Pri tem $N$ predstavlja število točk po koordinatni osi, torej je vseh točk v vektorju $u$ $N^2$. Pri tem opazim, da računski čas močno naraste, lastna vrednost pa se nad cca. $N=50$ malo spreminja. Presenetljivo je, da največja amplituda v prvi lastni obliki upada, vendar je konvergenca očitno počasnejša.
\begin{figure}[H]
    \centering
    \minifig{0.48}{./6.nal_slike/f_r100/cas_N}{Odvisnost računskega časa od št. točk na robu $N$}{}
    \minifig{0.48}{./6.nal_slike/f_r100/eigval_min_N}{Odvisnost najmanjše lastne vrednosti od št. točk na robu $N$}{}
\end{figure}
\fig{H}{0.48}{./6.nal_slike/f_r100/Xmax_N}{Odvisnost največje amplitude v prvi lastni obliki od št. točk na robu $N$}{}

Glede rezultatov je zanimivo, da ima ARPACK z večanjem/manjšanjem $ratio$ težave pri izračunu, dočim preostale metode ne. Prikažem le lastne oblike in lastne vrednosti za $ratio=100$ in $ratio=1/100$ in pri tem opazim, da večanje $ratio$ znižuje lastne vrednosti, zmanjševanje pa zvišuje, glede na homogeno opno ($ratio=1$).

Lastne oblike so prikazane za primer reševanja posplošenega problema. Pri sparse solverju so podobne, pri nesimetričnem pa sem imel malenkost več težav pri lastnih oblikah.

\fig{H}{0.8}{./6.nal_slike/f_r100/lastne}{Lastne oblike `f'-opne z $ratio = 100$}{}
\fig{H}{0.7}{./6.nal_slike/f_r100/eigval}{Lastne vrednosti, $ratio = 100$}{}

% Ratio = 1/100
\fig{H}{0.8}{./6.nal_slike/f_r0.01/lastne}{Lastne oblike `f'-opne z $ratio = 1/100$}{}
\fig{H}{0.7}{./6.nal_slike/f_r0.01/eigval}{Lastne vrednosti, $ratio = 1/100$}{}

\subsubsection{Kvadratna obtežba}

Zamislimo si še dva primera obtežbe, kjer je gostota na belem območju enaka $\rho_1=1$, na črnem pa $\rho_2=ratio$. Pri tem $ratio=\frac{\rho_2}{\rho_1}$ predstavlja razmerje gostot, enako kot prej.
\begin{figure}[H]
    \centering
    \subfig{0.3}{./6.nal_slike/opna_sr}{zgostitev na sredini}{}
    \subfig{0.3}{./6.nal_slike/opna_vogal}{zgostitev na vogalu}{}
    \caption{Porazdelitev gostote}
\end{figure}

%primerjava opne
Sledi primerjava za vse tri različne obtežbe opne in razmerja gostot $ratio$. Opazimo, da to ne vpliva na računski čas pri posplošenem načinu, kar je prikazano na spodnji sliki, vendar močno vpliva na sparse način.

Bolj zanimiva je odvisnost prve lastne vrednosti in največje amplitude prve lastne oblike, ki zgleda, da se spreminjajo zvezno. V splošnem se obe vrednosti zmanjšujeta z večanjem $ratio$.

\fig{H}{0.7}{./6.nal_slike/cas_ratio}{Odvisnost računskega časa od $ratio$}{}
\fig{H}{0.7}{./6.nal_slike/eigval_min_ratio}{Odvisnost med prvo lastno vrednostjo in $ratio$}{}
\fig{H}{0.7}{./6.nal_slike/Xmax_ratio}{Odvisnot med največjo amplitudo prve lastne oblike in $ratio$}{}

%primerjava last obl
Sledi še primerjava lastnih oblik pri različnih obtežbah, za kar uporabimo MAC (Modal Assurance Criterion) za statistično primerjavo (skaldnost/podobnost) dveh vektorjev lastnih oblik in je definiran kot
\begin{equation}
    \text{MAC}(\phi_l, \phi_k) = \frac{|\phi_l^{*T} \phi_k|^2}{(\phi_l^{*T} \phi_l)(\phi_k^{*T} \phi_k)},
\end{equation}
kjer $\phi_l^{*T}$ predstavlja kompleksno konjugirano in transponirano lastno obliko. Če sta $\phi_l$ in $\phi_k$ identični, bo MAC enak 1, če sta popolnoma različni pa 0.

Opazimo, da si lastne oblike bolj podobne pri nizkem $ratio$, kot pri visokem. Prav tako s približevanjem $ratio \rightarrow 1$ pričakovano narašča podobnost lastnih oblik, saj postajao opne tembolj homogene.

\begin{figure}[H]
    \centering
    \subfig{.45}{./6.nal_slike/MAC_opne/MAC_f_sr_r100}{f in sredina}{}
    \subfig{.45}{./6.nal_slike/MAC_opne/MAC_f_vogal_r100}{f in vogal}{}
    \subfig{.45}{./6.nal_slike/MAC_opne/MAC_sr_vogal_r100}{sredina in vogal}{}
    \caption{MAC za lastne oblike opne z različnimi oblikami obtežitev, $ratio=100$}
\end{figure}
\begin{figure}[H]
    \centering
    \subfig{.45}{./6.nal_slike/MAC_opne/MAC_f_sr_r10}{f in sredina}{}
    \subfig{.45}{./6.nal_slike/MAC_opne/MAC_f_vogal_r10}{f in vogal}{}
    \subfig{.45}{./6.nal_slike/MAC_opne/MAC_sr_vogal_r10}{sredina in vogal}{}
    \caption{MAC za lastne oblike opne z različnimi oblikami obtežitev, $ratio=10$}
\end{figure}
\begin{figure}[H]
    \centering
    \subfig{.45}{./6.nal_slike/MAC_opne/MAC_f_sr_r0.1}{f in sredina}{}
    \subfig{.45}{./6.nal_slike/MAC_opne/MAC_f_vogal_r0.1}{f in vogal}{}
    \subfig{.45}{./6.nal_slike/MAC_opne/MAC_sr_vogal_r0.1}{sredina in vogal}{}
    \caption{MAC za lastne oblike opne z različnimi oblikami obtežitev, $ratio=1/10$}
\end{figure}
\begin{figure}[H]
    \centering
    \subfig{.45}{./6.nal_slike/MAC_opne/MAC_f_sr_r0.01}{f in sredina}{}
    \subfig{.45}{./6.nal_slike/MAC_opne/MAC_f_vogal_r0.01}{f in vogal}{}
    \subfig{.45}{./6.nal_slike/MAC_opne/MAC_sr_vogal_r0.01}{sredina in vogal}{}
    \caption{MAC za lastne oblike opne z različnimi oblikami obtežitev, $ratio=1/100$}
\end{figure}

\section{Polkrožna opna}

Polkrožno opno obravnavamo v polarnem koordinatnem sistemu in dobimo podobno 5-diagonalno simetrično matriko $A$ kot prej, le z nekoliko drugačnimi koeficienti, ter matriko desnih strani $B$, v kolikor je opna nehomogena. Nehomogenost sem obravnaval na preprost način kot naraščanje/padanje gostote v radialni smeri. Robni pogoji so Dirichletovi $u=0$.

\begin{figure}[H]
    \centering
    \subfig{.48}{./6.nal_slike/zgosc_navzven/gostota}{zgoščevanje navzven}{}
    \subfig{.48}{./6.nal_slike/zgosc_noter/gostota}{zgoščevanje navznoter}{}
    \caption{Porazdelitev gostote med 1 in 10}
\end{figure}

Uporabil sem $N=N_r=N_{\phi}=50$ v radialni in obodni smeri. Pri uporabi sparse solverja iz \texttt{scipy.sparse.eigsh} za nehomogene opne sem ugotovil, da ima težave pri izračunu nizkih lastnih vrednosti, kar se odraža v slabi konvergenci in dolgih računskih časih. To sem preprosto odpravil z uporabo \href{https://docs.scipy.org/doc/scipy/tutorial/arpack.html} {shift-invert mode}, nato je sparse očitno hitrejši!. Pričakovano so pri nehomogenih opnah daljši računski časi.

Pri homogeni opni dobimo ustrezne lastne oblike, vendar sem pri sparse metodi imel težave, saj ni zaznala prve lastne oblike. Ugotovil sem, da je temu krivo dejstvo, da s parametrom $\sigma=0$ iščemo lastne vrednosti, ki so blizu 0, kar je pričakovano prva lastna oblika, vendar so v tem primeru prve lastne oblike negativne! Zato bi moral uporabiti  drugačen parameter. Taista težava je tudi pri inverzni potenčni metodi, ki je pri $\sigma=0$ konvergirala k višjim lastnim vrednostim.

Opazimo, da zgoščevanje navzven poudari prve lastne oblike in močno zniža lastne vrednosti, dočim za navznoter obratno. Na slikah lastnih oblik so prikazane absolutne vrednosti amplitud!

\fig{H}{1}{./6.nal_slike/polkrog/homog}{Lastne oblike homogene opne}{}

\fig{H}{1}{./6.nal_slike/zgosc_navzven/nehomog}{Lastne oblike pri zgoščevanju navzven}{}
\fig{H}{1}{./6.nal_slike/zgosc_noter/nehomog}{Lastne oblike pri zgoščevanju navznoter}{}

Opomba: Zaradi uporabi shift-invert metode pri sparse solverju, lastne vrednosti niso pravilno prikazane. Homogene so zamaknjene za 1, navzven zgoščevane za 4 in navznoter brez zamika. Ob upoštevanju tega so izračuni bolj skladni.
\fig{H}{.65}{./6.nal_slike/zgosc_navzven/eigval}{Lastne vrednosti homogene opne in zgoščevane navzven}{}
\fig{H}{.65}{./6.nal_slike/zgosc_noter/eigval}{Lastne vrednosti navznoter zgoščevane opne}{}

\fig{H}{.6}{./6.nal_slike/zgosc_navzven/čas}{računski čas homogene opne in zgoščevane navzven}{}
\fig{H}{.6}{./6.nal_slike/zgosc_noter/čas}{računski čas navznoter zgoščevane opne}{}